\documentclass[11pt,a4paper]{article}
\usepackage[utf8]{inputenc}
\usepackage[T1]{fontenc}
\usepackage[catalan]{babel}
\usepackage{lmodern}
\usepackage[sfdefault,lining]{FiraSans}
\usepackage[a4paper,top=2.3cm,bottom=2.2cm,left=1.8cm,right=1.8cm,headheight=1.3cm,headsep=0.4cm,footskip=1cm]{geometry}
\usepackage{amsmath,amssymb}
\usepackage{icomma}
\usepackage{graphicx}
\usepackage[table]{xcolor}
\usepackage{tikz}
\usepackage{fontawesome5}
\usepackage{ccicons}
\usepackage{fancyhdr}
\usepackage{lastpage}
\usepackage{tasks}
\usepackage{needspace}
\usepackage[most]{tcolorbox}
\usepackage{array}
\usepackage[hidelinks,pdftitle={Fitxa d'exercicis de la unitat 2: Divisibilitat. 1r d'ESO},pdfsubject={Matemàtiques 1r ESO},pdfkeywords={divisibilitat, múltiples, divisors, criteris de divisibilitat, nombres primers, factors primers, mcm, MCD}]{hyperref}

\definecolor{navy}{HTML}{1B2748}
\definecolor{accent}{HTML}{E8603A}
\definecolor{teal}{HTML}{1F9E90}
\definecolor{soft}{HTML}{EEF1F8}
\definecolor{muted}{HTML}{5D6685}

\setlength{\parindent}{0pt}
\setlength{\parskip}{3pt}
\settasks{label=\alph*), label-format=\bfseries\color{navy}, label-width=1.3em, item-indent=1.8em, column-sep=0.6em, after-item-skip=2pt, before-skip=3pt, after-skip=4pt}

% ---------- capçalera i peu ----------
\newcommand{\llicencia}{\href{https://creativecommons.org/licenses/by-nc-sa/4.0/deed.ca}{\ccbyncsa}\ \ Llicència Creative Commons Reconeixement-NoComercial-CompartirIgual 4.0 Internacional (CC BY-NC-SA 4.0)}
\pagestyle{fancy}
\fancyhf{}
\fancyhead[L]{\raisebox{-0.35\height}{\includegraphics[height=0.95cm]{logo-gva.png}}}
\fancyhead[C]{\small\color{muted}Fitxa d'exercicis · Unitat 2: Divisibilitat · 1r d'ESO}
\fancyhead[R]{\raisebox{-0.35\height}{\includegraphics[height=0.95cm]{logo-vermellar.png}}}
\fancyfoot[L]{\scriptsize\color{muted}\llicencia}
\fancyfoot[R]{\small\color{muted}\thepage}
\renewcommand{\headrulewidth}{0.4pt}
\renewcommand{\headrule}{\color{muted!40}\hrule height \headrulewidth}
\fancypagestyle{primera}{\fancyhf{}\renewcommand{\headrulewidth}{0pt}
  \fancyfoot[L]{\scriptsize\color{muted}\llicencia}
  \fancyfoot[R]{\small\color{muted}\thepage}}

% ---------- nivells i exercicis ----------
\newcommand{\estrelles}[1]{{\color{accent}\ifcase#1\or\faStar\faStar[regular]\faStar[regular]\or\faStar\faStar\faStar[regular]\or\faStar\faStar\faStar\fi}}
\newcounter{exer}
\newcommand{\ex}[1]{\par\needspace{4\baselineskip}\medskip\stepcounter{exer}\noindent{\color{navy}\bfseries\theexer.}\ \ {\small\estrelles{#1}}\ \ }

\newcommand{\seccio}[2]{\par\needspace{10\baselineskip}\bigskip
  \begin{tcolorbox}[colback=navy,colframe=navy,coltext=white,arc=2mm,boxrule=0pt,left=3mm,right=3mm,top=1.2mm,bottom=1.2mm]
  {\large\bfseries #1\quad #2}
  \end{tcolorbox}\vspace{-2pt}}

\newcolumntype{C}[1]{>{\centering\arraybackslash}p{#1}}

% quadre de recordatori (només a la fitxa d'exercicis)
\newcommand{\recorda}[1]{\par\smallskip\begin{tcolorbox}[colback=accent!7,colframe=accent!55,arc=2mm,boxrule=0.5pt,left=2.5mm,right=2.5mm,top=1mm,bottom=1mm]\small{\color{accent}\faLightbulb[regular]}\ \textbf{Recorda:}\ #1\end{tcolorbox}}
\newcommand{\res}[1]{#1}

\begin{document}
\thispagestyle{primera}

\begin{minipage}[c]{0.25\textwidth}\includegraphics[width=\linewidth]{logo-gva.png}\end{minipage}%
\begin{minipage}[c]{0.5\textwidth}\centering
  {\color{navy}\Large\bfseries Fitxa d'exercicis de la unitat 2:\\[2pt] Divisibilitat}\\[5pt]
  {\color{accent}\large\bfseries 1r d'ESO}
\end{minipage}%
\begin{minipage}[c]{0.25\textwidth}\hfill\includegraphics[width=0.85\linewidth]{logo-vermellar.png}\end{minipage}

\vspace{6pt}
{\color{navy}\hrule height 1.2pt}
\vspace{10pt}
Nom i cognoms: \rule{8.4cm}{0.4pt}\hfill Grup: \rule{1.6cm}{0.4pt}\hfill Data: \rule{2.2cm}{0.4pt}
\vspace{8pt}

\begin{tcolorbox}[colback=soft,colframe=soft,arc=2mm,left=3mm,right=3mm,top=1.5mm,bottom=1.5mm]
\small Exercicis graduats per nivells de dificultat:\quad \estrelles{1}\ bàsic \qquad \estrelles{2}\ mitjà \qquad \estrelles{3}\ avançat.\\[2pt] Resol-los al quadern mostrant tots els passos. La taula es pot completar directament en la fitxa. Fes els càlculs sense calculadora.
\end{tcolorbox}

% =====================================================
\seccio{1}{La relació de divisibilitat}

\ex{1} Fes la divisió i digues si el primer nombre és múltiple del segon:\begin{tasks}(4)
  \task $84$ i $7$
  \task $95$ i $8$
  \task $132$ i $12$
  \task $200$ i $15$
\end{tasks}\recorda{$a$ és \textbf{múltiple} de $b$ (i $b$ és \textbf{divisor} de $a$) quan la divisió $a:b$ és \textbf{exacta}: residu 0.}

\ex{1} Completa amb la paraula \emph{múltiple} o \emph{divisor}:\begin{tasks}(2)
  \task $35$ és \rule{2cm}{0.4pt} de $7$.
  \task $9$ és \rule{2cm}{0.4pt} de $63$.
  \task $100$ és \rule{2cm}{0.4pt} de $25$.
  \task $13$ és \rule{2cm}{0.4pt} de $52$.
\end{tasks}

\ex{1} A partir de la igualtat $8\cdot9=72$, escriu quatre frases amb les paraules \emph{múltiple} i \emph{divisor}.

\ex{2} Vertader o fals? Justifica-ho:\begin{tasks}(2)
  \task $48$ és divisible per $6$.
  \task $7$ és divisor de $77$.
  \task $120$ és múltiple de $9$.
  \task $0$ és múltiple de $15$.
  \task $1$ és múltiple de $3$.
\end{tasks}

\ex{2} Quin és el nombre més xicotet que cal sumar a $100$ perquè el resultat siga múltiple de $9$? I el més xicotet que cal restar-li?

\ex{3} En una divisió exacta, el divisor és $13$ i el quocient, $17$. Quin és el dividend? Amb el mateix divisor i quocient, quin és el dividend més gran que pot tindre una divisió \emph{no} exacta?

\ex{3} Troba tots els nombres de dues xifres que són múltiples de $7$ i acaben en $4$.


% =====================================================
\seccio{2}{Múltiples d'un nombre}

\ex{1} Escriu els deu primers múltiples (sense el 0) de:\begin{tasks}(3)
  \task $7$
  \task $11$
  \task $15$
\end{tasks}

\ex{1} Quins d'aquests nombres són múltiples de $6$? \[18\qquad 26\qquad 36\qquad 45\qquad 54\qquad 60\qquad 75\qquad 96\]

\ex{2} Escriu els múltiples de $13$ compresos entre $100$ i $170$.

\ex{2} Quants múltiples de $9$ hi ha entre $1$ i $200$? I entre $100$ i $200$?

\ex{2} Escriu els múltiples comuns de $6$ i de $9$ menors que $100$ (sense el 0). Quin és el més xicotet?

\ex{3} Quin és el múltiple de $17$ més pròxim a $300$?

\ex{3} Escriu el múltiple de $12$ més gran que tinga tres xifres i el més xicotet que en tinga quatre.


% =====================================================
\seccio{3}{Divisors d'un nombre}

\ex{1} Calcula tots els divisors de:\begin{tasks}(4)
  \task $15$
  \task $21$
  \task $28$
  \task $32$
\end{tasks}\recorda{Dividim entre $1, 2, 3\ldots$ Cada divisió exacta dona dos divisors (el divisor i el quocient). Parem quan el quocient és menor o igual que el divisor.}

\ex{1} Quins d'aquests nombres són divisors de $60$? \[4\qquad 7\qquad 9\qquad 12\qquad 15\qquad 25\qquad 30\]

\ex{2} Calcula tots els divisors de:\begin{tasks}(4)
  \task $48$
  \task $72$
  \task $90$
  \task $100$
\end{tasks}

\ex{2} Escriu els divisors comuns de $24$ i de $36$. Quin és el més gran?

\ex{2} En Pau té 36 cromos i els vol col·locar formant files, totes amb el mateix nombre de cromos. De quantes maneres ho pot fer?

\ex{3} Quin és el nombre menor que $50$ que té més divisors?

\ex{3} Troba un nombre que tinga exactament $4$ divisors i un altre que en tinga exactament $5$.


% =====================================================
\seccio{4}{Criteris de divisibilitat}

\ex{1} Completa la taula: marca amb una creu els nombres pels quals és divisible cada nombre.\par\smallskip{\renewcommand{\arraystretch}{1.5}\centering\begin{tabular}{|C{2.4cm}|C{1.5cm}|C{1.5cm}|C{1.5cm}|C{1.5cm}|C{1.5cm}|}\hline\rowcolor{soft}\textbf{Nombre} & $\boldsymbol{\div2}$ & $\boldsymbol{\div3}$ & $\boldsymbol{\div5}$ & $\boldsymbol{\div9}$ & $\boldsymbol{\div11}$\\\hline $3080$ &  &  &  &  & \\\hline $7425$ &  &  &  &  & \\\hline $1188$ &  &  &  &  & \\\hline $2543$ &  &  &  &  & \\\hline $6735$ &  &  &  &  & \\\hline \end{tabular}\par}\recorda{Per \textbf{2}: acaba en xifra parella. \ Per \textbf{3}: la suma de les xifres és múltiple de 3. \ Per \textbf{5}: acaba en 0 o 5. \ Per \textbf{9}: la suma de les xifres és múltiple de 9. \ Per \textbf{11}: la diferència entre la suma de les xifres de lloc senar i la de lloc parell és 0 o múltiple d'11.}

\ex{1} Escriu tres nombres de tres xifres que siguen divisibles per $3$ però no per $9$.

\ex{2} Quines xifres pots posar en lloc de $\square$ perquè $5\square4$ siga divisible per $3$? I per $9$?

\ex{2} Quina xifra cal posar en lloc de $\square$ perquè $38\square$ siga divisible per $2$ i per $3$ alhora?

\ex{2} Comprova, amb el criteri, si aquests nombres són divisibles per $11$:\begin{tasks}(4)
  \task $2849$
  \task $9375$
  \task $50\,182$
  \task $13\,575$
\end{tasks}

\ex{3} Troba el nombre més gran de tres xifres que siga divisible per $2$, per $3$ i per $5$ alhora.

\ex{3} Troba totes les xifres $a$ i $b$ perquè el nombre $7a2b$ siga divisible per $5$ i per $9$ alhora.


% =====================================================
\seccio{5}{Nombres primers i compostos}

\ex{1} Escriu els divisors de cada nombre i digues si és primer o compost:\begin{tasks}(3)
  \task $11$
  \task $15$
  \task $19$
  \task $27$
  \task $31$
  \task $49$
\end{tasks}\recorda{Un nombre és \textbf{primer} si només té dos divisors: l'1 i ell mateix. Si en té més, és \textbf{compost}. L'1 no és ni primer ni compost.}

\ex{1} Escriu tots els nombres primers compresos entre $60$ i $100$.

\ex{2} Esbrina si són primers dividint-los entre $2, 3, 5, 7, 11\ldots$ fins que el quocient siga menor que el divisor:\begin{tasks}(4)
  \task $87$
  \task $107$
  \task $133$
  \task $149$
\end{tasks}

\ex{2} Escriu cada nombre com a suma de dos nombres primers:\begin{tasks}(4)
  \task $16$
  \task $24$
  \task $36$
  \task $50$
\end{tasks}

\ex{2} Explica per què cap nombre acabat en $5$, més gran que $5$, pot ser primer. I per què cap nombre parell més gran que $2$ pot ser primer?

\ex{3} L'any $2027$ és un nombre primer? I el $2028$?

\ex{3} Troba tres nombres primers consecutius la suma dels quals siga també un nombre primer.


% =====================================================
\seccio{6}{Descomposició en factors primers}

\ex{1} Descompon en factors primers:\begin{tasks}(4)
  \task $36$
  \task $50$
  \task $72$
  \task $81$
\end{tasks}\recorda{Dividim entre el primer més xicotet possible (2, 3, 5, 7\ldots) fins a arribar a 1, i escrivim els factors repetits com a potència: $\ 360=2^{3}\cdot3^{2}\cdot5$.}

\ex{1} Descompon en factors primers:\begin{tasks}(4)
  \task $120$
  \task $126$
  \task $196$
  \task $225$
\end{tasks}

\ex{2} Descompon en factors primers:\begin{tasks}(4)
  \task $540$
  \task $784$
  \task $1260$
  \task $2310$
\end{tasks}

\ex{2} Quin nombre té cada descomposició?\begin{tasks}(4)
  \task $2^{4}\cdot3$
  \task $2\cdot3^{2}\cdot7$
  \task $2^{2}\cdot5^{3}$
  \task $3\cdot5\cdot11$
\end{tasks}

\ex{2} Completa les descomposicions:\begin{tasks}(3)
  \task $2^{\square}\cdot3\cdot5^{\square}=600$
  \task $2^{2}\cdot3^{\square}\cdot7=252$
  \task $2\cdot\square^{2}\cdot5=90$
\end{tasks}

\ex{3} Per quin nombre, el més xicotet possible, cal multiplicar $90$ perquè el resultat siga un quadrat perfecte?

\ex{3} Un nombre té la descomposició $2^{2}\cdot3\cdot5$. Quin nombre és? Escriu tots els seus divisors.


% =====================================================
\seccio{7}{Mínim comú múltiple}

\ex{1} Calcula el mcm escrivint la llista de múltiples:\begin{tasks}(4)
  \task $\text{mcm}(3,4)$
  \task $\text{mcm}(6,9)$
  \task $\text{mcm}(8,10)$
  \task $\text{mcm}(5,15)$
\end{tasks}

\ex{1} Calcula el mcm a partir de la descomposició en factors primers:\begin{tasks}(4)
  \task $\text{mcm}(12,20)$
  \task $\text{mcm}(18,30)$
  \task $\text{mcm}(24,36)$
  \task $\text{mcm}(14,21)$
\end{tasks}\recorda{\textbf{mcm}: factors \textbf{comuns i no comuns}, cadascun amb el \textbf{major} exponent.}

\ex{2} Calcula:\begin{tasks}(4)
  \task $\text{mcm}(40,60)$
  \task $\text{mcm}(36,54)$
  \task $\text{mcm}(45,75)$
  \task $\text{mcm}(28,42)$
\end{tasks}

\ex{2} Calcula:\begin{tasks}(3)
  \task $\text{mcm}(6,8,12)$
  \task $\text{mcm}(10,12,15)$
  \task $\text{mcm}(8,18,27)$
\end{tasks}

\ex{2} Calcula sense descompondre i explica com ho has fet:\begin{tasks}(4)
  \task $\text{mcm}(9,10)$
  \task $\text{mcm}(7,35)$
  \task $\text{mcm}(11,13)$
  \task $\text{mcm}(16,48)$
\end{tasks}

\ex{3} Calcula el mcm de $\ a=2^{3}\cdot3^{2}\cdot7\ $ i $\ b=2^{2}\cdot3^{3}\cdot5$.

\ex{3} El mcm de dos nombres és $60$ i un d'ells és $12$. Quins valors pot tindre l'altre?


% =====================================================
\seccio{8}{Problemes de mcm}

\ex{1} Un far s'encén cada $9$ segons i un altre, cada $12$. Si ara s'han encés tots dos alhora, d'ací a quant de temps tornaran a coincidir?\recorda{Si hi ha coses que es \textbf{repeteixen} i volem saber quan tornaran a \textbf{coincidir}, o busquem el nombre \textbf{més xicotet} que és múltiple de diversos, fem servir el \textbf{mcm}.}

\ex{1} La Laia entrena a futbol cada $3$ dies i va a natació cada $4$. Hui ha fet les dues coses. D'ací a quants dies les tornarà a fer el mateix dia?

\ex{2} Tres amics van al cinema: un cada $6$ dies, un altre cada $8$ i l'altre cada $10$. Si hui hi han coincidit, d'ací a quants dies tornaran a coincidir?

\ex{2} Quin és el nombre més xicotet de llapis que es pot empaquetar en capses de $12$, de $15$ o de $18$ sense que en sobre cap?

\ex{2} Dos autobusos ixen junts de la parada a les 7:30. Un hi torna cada $20$ minuts i l'altre, cada $25$. A quina hora tornaran a coincidir?

\ex{3} Si agrupem les galetes d'una caixa de $6$ en $6$, de $8$ en $8$ o de $10$ en $10$, sempre en sobren $3$. Quantes galetes hi ha com a mínim? I si sabem que n'hi ha entre $200$ i $300$?

\ex{3} Una roda dentada de $15$ dents engrana amb una altra de $24$ dents. Quantes voltes ha de fer cadascuna perquè tornen a estar en la posició inicial?


% =====================================================
\seccio{9}{Màxim comú divisor}

\ex{1} Calcula el MCD escrivint la llista de divisors:\begin{tasks}(4)
  \task $\text{MCD}(8,12)$
  \task $\text{MCD}(15,20)$
  \task $\text{MCD}(18,27)$
  \task $\text{MCD}(14,25)$
\end{tasks}

\ex{1} Calcula el MCD a partir de la descomposició en factors primers:\begin{tasks}(4)
  \task $\text{MCD}(24,36)$
  \task $\text{MCD}(30,45)$
  \task $\text{MCD}(42,56)$
  \task $\text{MCD}(60,90)$
\end{tasks}\recorda{\textbf{MCD}: \textbf{només} els factors \textbf{comuns}, cadascun amb el \textbf{menor} exponent. Si no n'hi ha cap de comú, el MCD és 1 (\emph{primers entre si}).}

\ex{2} Calcula:\begin{tasks}(4)
  \task $\text{MCD}(84,126)$
  \task $\text{MCD}(120,180)$
  \task $\text{MCD}(96,144)$
  \task $\text{MCD}(105,165)$
\end{tasks}

\ex{2} Calcula:\begin{tasks}(3)
  \task $\text{MCD}(12,18,30)$
  \task $\text{MCD}(24,40,56)$
  \task $\text{MCD}(45,75,105)$
\end{tasks}

\ex{2} Quines d'aquestes parelles són primers entre si?\begin{tasks}(4)
  \task $8$ i $21$
  \task $15$ i $25$
  \task $14$ i $27$
  \task $22$ i $33$
\end{tasks}

\ex{3} Calcula el MCD i el mcm de $36$ i $60$, i comprova que $\text{MCD}\cdot\text{mcm}=36\cdot60$.

\ex{3} El MCD de dos nombres és $8$ i el seu mcm, $96$. Si un dels nombres és $24$, quin és l'altre?


% =====================================================
\seccio{10}{Problemes de MCD}

\ex{1} Tenim $30$ pomes i $45$ taronges i volem fer el màxim nombre de lots iguals, sense que sobre cap fruita. Quants lots podem fer? Què portarà cada lot?\recorda{Si volem \textbf{repartir} o \textbf{tallar} en parts iguals, \textbf{sense que sobre res}, i les parts han de ser \textbf{el més grans possible} (o el màxim nombre de grups), fem servir el \textbf{MCD}.}

\ex{1} Volem tallar dues cordes de $36$ m i $48$ m en trossos iguals, el més llargs possible. Quant farà cada tros? Quants trossos tindrem?

\ex{2} Volem enrajolar un terra de $360$ cm per $280$ cm amb rajoles quadrades, el més grans possible, sense tallar-ne cap. Quant farà el costat de cada rajola? Quantes rajoles calen?

\ex{2} Una llibreria té $120$ llibres de contes, $84$ de poesia i $60$ còmics. Vol fer paquets amb el mateix nombre de llibres, tots d'un sol tipus i el més grans possible. Quants llibres hi haurà en cada paquet? Quants paquets farà?

\ex{2} Digues si cal calcular el mcm o el MCD i resol:\begin{tasks}(1)
  \task Un rellotge sona cada $15$ minuts i un altre, cada $25$. Si ara han sonat junts, quan tornaran a coincidir?
  \task Volem repartir $48$ bolígrafs i $64$ llapis en el màxim nombre d'estoigs iguals, sense que en sobre cap.
\end{tasks}

\ex{3} Una mestra reparteix $75$ cromos i $105$ adhesius entre l'alumnat, de manera que tots reben la mateixa quantitat de cada cosa. Li sobren $3$ cromos i $9$ adhesius. Quants alumnes hi pot haver, com a màxim? Quants cromos i adhesius rep cadascú?

\ex{3} Volem plantar arbres al voltant d'un jardí rectangular de $72$ m per $48$ m, a la mateixa distància entre ells, la més gran possible i amb un arbre a cada cantonada. A quina distància els plantarem? Quants arbres calen?

\vfill
\begin{center}\color{muted}\small\faCheckCircle[regular]\ \ Ànim! Si et costa algun exercici, repassa els apunts interactius de la unitat.\end{center}

\end{document}
